% =============================================================================
% 04_discussion.tex -- Discussion of physical implications and constraints
% Fig.~\ref{fig:kkg-cdm-density-3d}: code/calc-07.py -> manuscript/figures/figure_7.pdf
% Fig.~\ref{fig:kkg-tov-mr}: code/calc-10.py -> manuscript/figures/figure_10.pdf
% =============================================================================

\section{Discussion}
\label{sec:discussion}

The Phase-Locked Klein--Gordon (PLKG) model developed in this work bridges two traditionally separate domains: non-linear synchronization dynamics and cosmological scalar-field theories. In this section, we explore the physical implications of this coupling, focusing on three key areas: the fuzzy dark matter regime, the emergence of synchronization pressure, and potential observational constraints.

\subsection{Fuzzy Dark Matter and Synchronization}
\label{subsec:fdm_sync}

In the ultra-light mass regime $m \sim 10^{-22}\,\text{eV}$, Klein-Gordon fields behave as \emph{fuzzy dark matter} (FDM) \cite{Hu2000,Marsh2016}, where the de~Broglie wavelength becomes galactic-scale and quantum pressure suppresses small-scale structure formation. The introduction of Kuramoto coupling in this regime adds a \emph{classical coherence mechanism} that may fundamentally alter the dynamics of structure formation.

\subsubsection{Synchronization Pressure}
The interaction contribution to the comoving energy density $\rho_{\mathrm{tot}}$ matches the explicit cosine form in Sec.~\ref{subsec:emt},
\begin{equation}
\rho_{\text{int}} = \frac{K}{N}\sum_{j,k=1}^N \rho_j\rho_k\cos(\theta_k-\theta_j) = -V_{\mathrm{int}}\,,
\end{equation}
does not by itself define a stiff fluid: viewed in isolation, the bare potential contribution is vacuum-like with $w_{V_{\text{int}}}=-1$. However, in the synchronized dynamical regime the full PLKG stress tensor can be dominated by coherent phase-kinetic stiffness, yielding an effective equation of state $w_{\text{sync}} \approx 1$ for the conserved synchronized sector. This effective ``synchronization pressure'' acts to:
\begin{itemize}
    \item \textbf{Suppress fragmentation:} Fully synchronized populations $(\theta_k = \theta_j)$ resist gravitational collapse on scales smaller than the synchronization length
    \begin{equation}
    \lambda_{\text{sync}} \sim \frac{1}{\sqrt{K}},
    \label{eq:lambda_sync}
    \end{equation}
    where $K$ is the Kuramoto coupling constant. Since $V_{\text{int}} \sim K\rho^2$ and $[\rho]=M$ in natural units, one has $[K]=M^2$, so the combination $1/\sqrt{K}$ has the correct length dimension $M^{-1}$. Thus the synchronization scale is set directly by the coupling scale, without requiring an additional factor of $m$ for dimensional consistency.
    \item \textbf{Modify halo profiles:} The interplay between quantum pressure (from $m$) and synchronization pressure (from $K$) may produce core-halo structures distinct from standard FDM predictions.
    \item \textbf{Enhance coherence:} Unlike pure FDM, where coherence arises solely from quantum interference, the Kuramoto coupling introduces a \emph{dynamical attractor} toward phase alignment, potentially stabilizing solitonic cores.
\end{itemize}

\subsubsection{Illustrative cusp versus core morphology}
\label{subsubsec:cusp-core-cartoon}

At the level of a single spatial dimension and a schematic time axis, standard CDM tends to deepen a central density cusp as material concentrates at small radii.
By contrast, a PLKG-motivated stiff regime ($w_{\text{sync}}\approx 1$ when a synchronization order parameter $R$ is near unity) suggests an additional effective ``synchronization pressure'' that can flatten the inner profile into a core-like morphology.
Figure~\ref{fig:kkg-cdm-density-3d} compares two closed-form surfaces $\rho(x,t)$ (a narrowing Gaussian for CDM and a stiffening super-Gaussian tied to $R(t)$ for PLKG) to visualize this qualitative contrast.
It is \emph{not} a numerical relativity or N-body realization.
The figure is generated by \texttt{code/calc-07.py} (\texttt{--prd-width column} or \texttt{--prd-width full} sets the canvas geometry before \verb|\includegraphics| scaling).

\begin{figure*}[t]
    \centering
    \includegraphics[width=\textwidth]{figures/figure_7.pdf}
    \caption{Schematic three-dimensional comparison of probability-of-presence profiles $\rho(x,t)$.
    \textbf{Left (viridis):} CDM-like cusp: a normalized Gaussian whose width shrinks with time, concentrating density at $x=0$.
    \textbf{Right (magma):} PLKG-like core: a super-Gaussian whose effective exponent grows with a proxy order parameter $R(t)\to 1$, producing a progressively flatter central plateau and stiffer shoulders (illustrative synchronization pressure).
    Neither surface is observational data; the left panel caricatures the cuspy endmember of DM-only halo models, while the right panel caricatures a centrally flattened profile of the kind often inferred from dwarf-galaxy rotation curves (see text and \protect\cite{Bullock2017}).
    Generated by \texttt{code/calc-07.py}.}
    \label{fig:kkg-cdm-density-3d}
\end{figure*}

The \textbf{resulting morphology} in Fig.~\ref{fig:kkg-cdm-density-3d} can be read as follows.
The CDM panel shows a sharpening central peak as time increases: this is a cartoon for the \emph{cusp} tendency of dissipationless clustering in pure dark-matter-only models.
The PLKG panel shows the opposite qualitative trend for the chosen proxy: the central region becomes \emph{flatter} while the profile steepens farther out, a cartoon for a \emph{core} supported by additional effective pressure when synchronization becomes efficient.
Neither surface is fitted to data or obtained from the full covariant PLKG system in three spatial dimensions with gravity.

Turning to \textbf{observations}, neither panel should be mistaken for a measured $\rho(r)$.
However, at the level of \emph{which qualitative shape is closer to widely discussed rotation-curve systematics}, cold collisionless DM-only simulations typically produce \emph{cuspy} inner halos, while many analyses of dwarf and low-surface-brightness galaxies favor \emph{centrally flattened} circular-speed curves compared to naive cuspy fits \cite{Bullock2017}.
Thus the \textbf{right} (PLKG-motivated core) panel is closer to that \emph{phenomenological} ``core'' direction in some disk galaxies, while the \textbf{left} (CDM cusp) panel is closer to the \emph{dark-matter-only} simulation endmember.
Those cores are not unique to PLKG: baryonic feedback, self-interacting DM, and fuzzy-DM wave interference can also flatten inner profiles \cite{Bullock2017}, so identifying a PLKG contribution requires predictions beyond this schematic.

\subsubsection{Microscopic $|\Phi|^2$ stability under shared noise (toy model)}
\label{subsubsec:phi2-toy}

To complement the macroscopic cusp--core cartoon of Fig.~\ref{fig:kkg-cdm-density-3d}, we also record a minimal one-dimensional toy in which $\Phi(x,t)=\sum_j \rho_j(x)\,e^{i\theta_j(t)}$ is built from fixed Gaussian envelopes $\rho_j(x)$ on $x\in[-10,10]$ with $N=40$ modes.
Phases follow an Euler--Maruyama update with identical white-noise draws for two runs: an FDM-like sector with $K=0$ (no Kuramoto torque) and a PLKG sector with $K=5$, intrinsic frequencies centered on $m=1$, noise strength $\eta=0.3$, and $\Delta t=0.05$ for $200$ steps (\texttt{code/calc-08.py}).
Figure~\ref{fig:kkg-fdm-phi-stability} compares $|\Phi|^2$ in $(x,t)$ and the time series of the spatial maximum $M(t)=\max_x|\Phi|^2$.
The lower panel reports the temporal variance of $M(t)$ after a short burn-in; in this schematic, the synchronized run exhibits substantially smaller variance, illustrating how coupling can suppress noise-driven modulation of the peak density relative to the incoherent case.
This demonstration is not a substitute for a full quantum FDM or covariant PLKG simulation.

\begin{figure*}[t]
    \centering
    \includegraphics[width=\textwidth]{figures/figure_8.pdf}
    \caption{Toy microscopic comparison of $|\Phi(x,t)|^2$ for multi-mode Gaussian packets with shared phase noise.
    (A) $K=0$: phases decohere independently (FDM-like sector).
    (B) PLKG-like ($K>0$): phase-locked coupling + same noise.
    (C) $M(t)=\max_x|\Phi|^2$ and temporal variance after burn-in (defaults in \texttt{code/calc-08.py}).
    Neither heat map is observational data.
    Generated by \texttt{code/calc-08.py}.}
    \label{fig:kkg-fdm-phi-stability}
\end{figure*}

\textbf{Relation to observations.}
Neither panel of Fig.~\ref{fig:kkg-fdm-phi-stability} is a measurement of $|\Phi|^2$ in a galaxy.
The $K=0$ trajectory is only an \emph{FDM-like} classical caricature in which independent noisy phases shred interference structure; it is \emph{not} the output of a Schr\"odinger--Poisson solver for ultra-light bosons, where quantum coherence and wave pressure are essential \cite{Hu2000,Marsh2016}.
The $K>0$ trajectory is likewise a classical Kuramoto-on-modes cartoon, not a covariant PLKG field simulation coupled to gravity.
At the level of \emph{phenomenology}, however, inferred inner mass distributions in dwarfs and low-surface-brightness systems are usually discussed as smooth, slowly varying potentials on orbital timescales rather than violently fluctuating central densities \cite{Bullock2017}.
In that loose sense, the synchronized run---with a much smaller temporal variance of $\max_x|\Phi|^2$ under identical noise draws---is closer to the \emph{direction} of ``stable, coherent'' inner structure that such rotation-curve analyses motivate, whereas the incoherent run illustrates a noise-dominated fragmentation of the envelope that should not be over-identified with full FDM predictions.
Quantitative contact with data still requires solving the appropriate linear and nonlinear field equations (including baryons) and comparing to survey forward models.

\subsubsection{Desynchronization and Effective Pressurelessness}
\label{subsubsec:desync}
In contrast, a \emph{desynchronized population} with uniformly distributed phases $\theta_j \in [0, 2\pi)$ yields
\begin{equation}
\langle \cos(\theta_j - \theta_k) \rangle = 0,
\end{equation}
causing the interaction energy to average out. In this regime, the system behaves as \emph{pressureless dust} ($w \approx 0$), mimicking cold dark matter (CDM). This suggests a \emph{dynamical transition} from stiff-fluid-like behavior at early times (when synchronization is strong) to CDM-like behavior at late times (when cosmic expansion or gravitational perturbations destroy coherence).

\subsection{Cosmological Implications}
\label{subsec:cosmo_implications}

\subsubsection{Early Universe Constraints}
\label{subsubsec:early-universe}

Throughout this subsection the Kuramoto coupling is quoted in units of $\mathrm{eV}^2$, consistent with mass dimension $[K]=M^2$ for $V_{\mathrm{int}}\sim K\rho^2$ in natural units.

\paragraph{(a) Synchronized versus desynchronized limits.}
In a spatially flat FLRW background, a perfect fluid with equation-of-state parameter $w\equiv p/\rho$ satisfies $\rho\propto a^{-3(1+w)}$ when the fluid does not exchange energy with other sectors.
Radiation has $w=1/3$ and $\rho_{\mathrm{rad}}\propto a^{-4}\propto T^4$ under adiabatic expansion.
Pressureless matter has $w=0$ and $\rho_{m}\propto a^{-3}\propto T^3$.
A stiff source with $w=1$ scales as $\rho_{\mathrm{stiff}}\propto a^{-6}\propto T^6$.
The effective synchronized sector discussed in Sec.~\ref{subsec:fdm_sync} can behave approximately as a stiff fluid, $w_{\mathrm{sync}}\approx 1$, when phase coherence is efficient and the associated energy density is dominated by rapidly varying, coherently locked phase degrees of freedom.
By contrast, a uniformly desynchronized population averages the interaction stress so that the coherent stiff channel shuts off and the sector behaves as pressureless dust, $w\approx 0$, as summarized in Sec.~\ref{subsubsec:desync}.
These two limits are not interchangeable with a radiation-like fluid: stiff matter is not ``extra radiation,'' and constraints phrased for $w\approx 1/3$ (for example, bounds on the effective number of neutrino species) do not map onto a putative $w=1$ component without re-deriving the Boltzmann hierarchy for that equation of state.

\paragraph{(b) BBN constraints on a stiff PLKG component.}
Big-bang nucleosynthesis fixes the expansion rate at temperatures $T\sim 0.05$--$1\,\mathrm{MeV}$ through the standard timeline for neutron freeze-out and nuclear reaction rates \cite{Cyburt2016}.
Additional energy density in a species with $w\neq 1/3$ enters the Friedmann equation directly through $\rho_{\mathrm{tot}}$ and is not usefully summarized as a shift in $N_{\mathrm{eff}}$, which is defined assuming free-streaming radiation with $w=1/3$.

Suppose the PLKG sector could be modeled as stiff from matter--radiation equality all the way back to BBN, and that it contributes a fraction $\Omega_{\mathrm{PLKG},0}\sim\Omega_{\mathrm{dm},0}\sim 0.3$ of critical density today while remaining subdominant to photons in the budget today, $\Omega_{\mathrm{rad},0}\sim\mathcal{O}(10^{-4})$.
Then, comparing the same sector to photons plus neutrinos modeled collectively as radiation with $\rho_{\mathrm{rad}}\propto a^{-4}$,
\begin{equation}
\frac{\rho_{\mathrm{PLKG}}}{\rho_{\mathrm{rad}}}\Big|_{\mathrm{BBN}}
=\frac{\Omega_{\mathrm{PLKG},0}}{\Omega_{\mathrm{rad},0}}
\left(\frac{a_0}{a_{\mathrm{BBN}}}\right)^{2}
\sim\frac{0.3}{10^{-4}}\left(\frac{T_{\mathrm{BBN}}}{T_0}\right)^{2}
\sim\mathcal{O}(10^{23}),
\label{eq:stiff-bbn-disaster}
\end{equation}
where $T_{\mathrm{BBN}}/T_0\sim 10^{10}$ for $T_{\mathrm{BBN}}\sim\mathcal{O}(1)\,\mathrm{MeV}$ and $T_0\sim 10^{-10}\,\mathrm{MeV}$.
Such a huge ratio would drive the Hubble rate at BBN far above the baseline used in standard primordial abundance fits and is therefore excluded at the level of the scaling argument alone \cite{Cyburt2016}.
In particular, a stiff dark-matter candidate cannot simply track $\Omega_{\mathrm{dm}}$ today while retaining $w=1$ back to $T\sim 1\,\mathrm{MeV}$ unless its abundance was utterly negligible during BBN.

\paragraph{(c) BBN constraints on a pressureless PLKG component.}
If instead the PLKG fields are desynchronized at $T\sim 1\,\mathrm{MeV}$, then $w\approx 0$ and $\rho_{\mathrm{PLKG}}\propto a^{-3}$, the same scaling as non-relativistic matter.
Relative to radiation, $\rho_{\mathrm{PLKG}}/\rho_{\mathrm{rad}}\propto a\propto T^{-1}$, so a subdominant cold relic redshifts away faster than photons and becomes more negligible at earlier times.
The primary BBN impact of a cold relic that contributes only a small fraction of $\rho_{\mathrm{tot}}$ is suppressed compared to the stiff case in paragraph~(b); standard analyses then focus on the relic's allowed energy density at equality and on structure-formation constraints rather than on a large fractional change to the expansion rate during neutron freeze-out \cite{Cyburt2016}.

\paragraph{(d) Transition epoch from Hubble-limited synchronization.}
The overdamped Kuramoto limit of the main text implies that cosmic expansion competes with mean-field locking through the dimensionless ratio $K/(3H)$ in the phase evolution.
Appendix~\ref{app:stability} packages the same physics in the criterion that cosmologically efficient synchronization within one Hubble time requires $K\gtrsim K_{\mathrm{cosmo}}\sim 3H(z)N$ \eqref{eq:sync_threshold}.
Equivalently, for
\begin{equation}
3H(z)N \gg K,
\label{eq:hubble-wins}
\end{equation}
Hubble drag prevents rapid alignment of the phases, the Kuramoto order parameter is driven toward small values, and the discussion of Sec.~\ref{subsubsec:desync} indicates that the stiff channel is shut off in favor of an averaged, pressureless-like equation of state.
Define a characteristic redshift $z_{\mathrm{trans}}$ by
\begin{equation}
3H(z_{\mathrm{trans}})N \sim K.
\label{eq:ztrans-def}
\end{equation}
During radiation domination, ignoring changes in the effective number of relativistic species,
\begin{equation}
H(z)\approx H_0\sqrt{\Omega_{r,0}}\,(1+z)^2,
\label{eq:H-rad-approx}
\end{equation}
so that
\begin{equation}
(1+z_{\mathrm{trans}})^2\sim\frac{K}{3N H_0\sqrt{\Omega_{r,0}}}.
\label{eq:ztrans-scaling}
\end{equation}
Whenever $z_{\mathrm{BBN}}\gg z_{\mathrm{trans}}$, the universe at BBN sits deeply in the regime \eqref{eq:hubble-wins} and the PLKG sector is pressureless for the purpose of the scaling in paragraph~(c), evading the catastrophe \eqref{eq:stiff-bbn-disaster}.
The complementary logical possibility---that almost all of today's PLKG abundance was produced after BBN---requires a separate entropy-production or coupling history that the present mean-field construction does not specify; we do not pursue it here.

\paragraph{(e) CMB-era constraints and dimensional bounds on $K$.}
For $z\sim 10^{3}$--$10^{4}$, the same comparison $3H(z)N$ versus $K$ determines whether a stiff synchronized channel can reactivate after BBN.
If the sector remains pressureless and clusters like cold dark matter, compatibility with the cosmic microwave background is captured at leading order by the usual bounds on non-baryonic matter density and background equation of state, e.g.\ $\Omega_{\mathrm{cdm}}h^2$ and constraints on dynamical dark energy, rather than by $N_{\mathrm{eff}}$ \cite{Planck2018}.

Independently of cosmology, linear stability of the homogeneous amplitude requires $m_{\mathrm{eff}}^2=m^2-K-\dot{\bar{\theta}}^2>0$ in the notation of Appendix~\ref{app:stability} \eqref{eq:meff}.
Thus $K$ must obey $K\lesssim m^2$ up to $\mathcal{O}(1)$ factors set by the background phase velocity.
For fuzzy-dark-matter masses $m\sim 10^{-22}\,\mathrm{eV}$, a conservative stability-oriented scale is
\begin{equation}
K\lesssim m^2\sim 10^{-44}\,\mathrm{eV}^2,
\label{eq:K-stability-scale}
\end{equation}
quoted with units of $\mathrm{eV}^2$ as required by $[K]=2$.
Any scenario in which a stiff component with $w\approx 1$ survives until recombination would modify the acoustic history and early integrated Sachs--Wolfe sources through the altered $H(z)$; extracting tight bounds on that sector requires explicitly solving the coupled perturbation equations for the PLKG variables and comparing to Planck likelihoods \cite{Planck2018}, which we leave to future work.

\subsubsection{Structure Formation}
The PLKG model predicts a \emph{scale-dependent suppression} of the matter power spectrum:
\begin{itemize}
    \item \textbf{Small scales ($k \gg k_{\text{sync}}$):} Synchronization pressure suppresses power beyond the standard FDM cutoff, potentially resolving the ``missing satellites'' problem more effectively than pure FDM.
    \item \textbf{Large scales ($k \ll k_{\text{sync}}$):} The model reduces to standard FDM or CDM, preserving agreement with large-scale structure observations.
\end{itemize}

\subsection{Compact bosonic stars and a schematic TOV comparison}
\label{subsec:tov-compact}

Gravitational-wave catalogs have highlighted a phenomenological ``mass gap'' around roughly $2.5$--$3\,M_\odot$, where a handful of primary components sit uncomfortably between canonical neutron-star upper bounds inferred from nuclear equations of state and the lightest stellar-mass black holes.
Within the PLKG program, a synchronized scalar sector can supply an additional stiffening channel beyond a baseline self-gravitating Bose condensate: collective phase alignment carries an energy cost that can be modeled schematically as a positive-definite correction to the barotropic pressure, $P_{\mathrm{tot}}=P_{\mathrm{poly}}+P_{\mathrm{sync}}$ with $P_{\mathrm{sync}}\propto K\rho^2$ in the same spirit as the phase-stiffening discussion above.
Figure~\ref{fig:kkg-tov-mr} shows a \emph{toy} integration of the Tolman--Oppenheimer--Volkoff equations in geometric units ($G=c=1$) for a fixed polytrope $P_{\mathrm{poly}}\propto\rho^{\beta}$ plus the $K$-dependent term; it is not a microphysical derivation from the covariant PLKG stress tensor, nor is it calibrated to NICER or GW data.
Canvas sizing follows the same \texttt{--prd-width column|full} convention as the other \texttt{code/calc-*.py} figure drivers.
Qualitatively, increasing $K$ raises the maximum stable mass and shifts the mass--radius knee toward larger radii at fixed mass compared to the $K=0$ baseline, illustrating how synchronization-mediated stiffness could push bosonic configurations into the intermediate-mass window without invoking primordial black holes alone.
Relative to standard Bose--Einstein-condensate star models with comparable polytropic cores, the Kuramoto coupling here acts as a \emph{collective} reinforcement: phase locking distributes the burden of supporting self-gravity across the locked population, postponing the onset of the unstable branch in this barotropic closure.
Finally, a radially stratified PLKG star naturally suggests a sharp \emph{transition layer} between a synchronized interior (large effective stiffness) and a desynchronized crust (pressureless-like averaging of random phases).
Sudden angular-momentum exchange during crust--core spin-up events could then be reinterpreted, at least at a cartoon level, as pulsar glitches triggered when the crust intermittently slips relative to a more rigidly synchronized core, i.e.\ as a finite-amplitude rearrangement across that internal coherence transition rather than solely as a crustquake in the traditional nuclear sense.

\begin{figure*}[t]
    \centering
    \includegraphics[width=\textwidth]{figures/figure_10.pdf}
    \caption{Toy Tolman--Oppenheimer--Volkoff mass--radius sequences comparing a baseline barotropic star ($P_{\mathrm{poly}}=\kappa\rho^{\beta}$) to the same model augmented by $P_{\mathrm{sync}}=\alpha_{\mathrm{sync}}K\rho^2$, integrated outward from a dense core until the surface pressure threshold is reached ($G=c=1$; boson mass scale fixed by the chosen polytrope coefficients, not fitted to observations).
    (a) Mass--radius curves for several $K$; markers indicate the numerically identified turnover on the stable branch.
    (b) Peak mass $M_{\max}$ versus $K$.
    Generated by \texttt{code/calc-10.py}.}
    \label{fig:kkg-tov-mr}
\end{figure*}

\subsection{Observational Signatures}
\label{subsec:obs_signatures}

\textbf{Halo mass function:}
The synchronization length $\lambda_{\text{sync}}$ introduces
a characteristic scale below which halo formation is suppressed...

\textbf{Rotation curves:}
In dwarf galaxies ($M \sim 10^8 M_\odot$), synchronization
pressure modifies the density profile...

\textbf{Lyman-$\alpha$ forest:}
Current constraints on the matter power spectrum at
$k \sim 1\,h\,\text{Mpc}^{-1}$ imply...

\textbf{CMB and BBN:}
The coupling constant $K$ must satisfy $K/H_{\text{BBN}} < ...$
to avoid disrupting nucleosynthesis.


\subsection{Comparison with Related Models}
\label{subsec:comparison}

\subsubsection{Standard Kuramoto Model}
The classical Kuramoto model on a network assumes \emph{fixed spatial positions} and instantaneous interactions. In contrast, the PLKG model:
\begin{itemize}
    \item Treats oscillators as \emph{dynamical scalar fields} propagating on curved spacetime.
    \item Incorporates \emph{retardation effects} from finite light-speed.
    \item Couples synchronization to \emph{gravitational dynamics} via the stress-energy tensor.
\end{itemize}

\subsubsection{Scalar-Field Dark Matter}
Compared to standard scalar-field dark matter models (e.g., axion-like particles, ultralight dark matter), the PLKG model introduces:
\begin{itemize}
    \item \textbf{Self-interaction:} The Kuramoto term acts as a \emph{non-local, phase-dependent interaction}, distinct from quartic self-interactions $\lambda \phi^4$.
    \item \textbf{Emergent coherence:} Synchronization arises \emph{dynamically} rather than being imposed as an initial condition.
\end{itemize}

\subsection{Limitations and Future Directions}
\label{subsec:limitations}

\subsubsection{Model Limitations}
\begin{enumerate}
    \item \textbf{Mean-field approximation:} The Kuramoto coupling assumes all-to-all interactions, which may not hold in realistic cosmological settings where causality limits interaction ranges.
    \item \textbf{Classical treatment:} Quantum corrections (e.g., loop effects, entanglement) are neglected. A full quantum field theory treatment may reveal new phenomena.
    \item \textbf{Homogeneity assumption:} The effective stiff-fluid result $w_{\text{sync}} \approx 1$ assumes a homogeneous, phase-locked regime in which the conserved synchronization energy is dominated by coherent phase motion. Inhomogeneous configurations (e.g., near black holes) require numerical simulations.
\end{enumerate}

\subsubsection{Future Work}
\begin{itemize}
    \item \textbf{Numerical simulations:} Solve the coupled Einstein--PLKG equations in cosmological and astrophysical contexts (e.g., galaxy formation, black hole accretion).
    \item \textbf{Observational tests:} Compute predictions for:
    \begin{itemize}
        \item Halo mass functions and concentration-mass relations.
        \item Lyman-$\alpha$ forest constraints on small-scale power.
        \item Gravitational wave signatures from synchronized scalar-field oscillations.
    \end{itemize}
    \item \textbf{Generalizations:}
    \begin{itemize}
        \item Extend to \emph{non-uniform coupling} $K_{jk}(x)$ (e.g., distance-dependent interactions).
        \item Include \emph{frequency disorder} $\omega_j \neq \omega_k$ to model realistic astrophysical populations.
        \item Couple to \emph{other fields} (e.g., electromagnetic fields, neutrinos).
    \end{itemize}
\end{itemize}

\subsection{Concluding Remarks}
\label{subsec:concluding-remarks}
The phase-locked Klein--Gordon (PLKG) framework represents a novel synthesis of synchronization theory and relativistic field theory. By introducing a dynamical mechanism for phase coherence in scalar-field dark matter, it offers new avenues for addressing long-standing puzzles in cosmology and astrophysics. While significant theoretical and observational work remains, the framework developed here provides a rigorous foundation for exploring the interplay between collective dynamics and gravity.
